% ── Section VI: Conclusion ───────────────────────────────────────────
% Length  : ~0.5 page
% Structure:
%   1. Summary     — what was done and the main result (2–3 sentences)
%   2. Limitations — honest, 2–3 bullets
%   3. Future Work — forward pointer to Paper 3 (edge deployment)
% ─────────────────────────────────────────────────────────────────────

\section{Conclusion}
\label{sec:conclusion}

% [Summary]
% TODO: restate method and quantitative improvement
%   "We proposed a context-aware adaptive perception pipeline for
%    motorcycle-dominant mixed traffic, comprising three novel
%    contributions: [C1, C2, C3].
%    On our self-collected Vietnamese urban dataset, the full pipeline
%    achieves XX\% MOTA and reduces false-positive alert rate by XX\%
%    compared to the standard ByteTrack baseline, while maintaining
%    real-time performance at XX~FPS."

% [Limitations]
% TODO: be honest — reviewers respect self-awareness
% Suggested limitations:
%   - Dataset is single-city; generalizability to other SEA cities untested
%   - Context estimator uses heuristic thresholds; learned adaptation is future work
%   - Evaluation on embedded hardware deferred to Paper 3

% [Future Work — Paper 3]
% TODO: forward pointer
%   "In future work, we plan to deploy this pipeline on the
%    NVIDIA Jetson Orin NX platform as part of a full-stack edge
%    system integrating adaptive perception with multi-agent
%    signal control, targeting real-time performance within
%    the constraints of embedded hardware."
